% !TeX encoding = UTF-8
\documentclass[variant=docente,cover=institutional,mode=screen]{../cls/ua-engineering-v6-article}
% !TeX encoding = UTF-8
\UASetUniversity{Universidad Austral}
\UASetFaculty{Facultad de Ingeniería}
\UASetDepartment{Departamento de Computación}
\UASetCampus{Pilar}
\UASetCareer{Ingeniería Informática}
\UASetCommission{Comisión A}
\UASetProfessor{Equipo docente}
\UASetSemester{Segundo cuatrimestre 2026}
\UASetCourse{Sistemas Distribuidos}
\UASetDate{2026}


\UASetClassNumber{02}
\UASetClassTitle{RPC, timeouts, retries e idempotencia}
\UASetDocKind{Guía docente}

\title{Clase 02 --- Guía docente}
\author{\UAProfessor}
\date{\UADate}

\begin{document}
\maketitle

\section{Objetivo profundo}

El objetivo de esta clase es que el estudiante entienda que:

\begin{itemize}
\item RPC no es transparente
\item la comunicación distribuida no es inmediata ni atómica
\item una llamada atraviesa capas que pueden observar y responder distinto ante la falla
\item timeout introduce incertidumbre inevitable
\item retry introduce duplicación estructural
\item idempotencia es la solución clave
\end{itemize}

\section{Resultado esperado}

El alumno debe poder explicar:

\begin{itemize}
\item por qué RPC rompe la analogía con llamadas locales
\item cómo reconstruir una comunicación remota como línea de tiempo
\item qué capas intervienen y qué puede hacer cada una ante pérdida, demora o saturación
\item por qué retry no es opcional
\item por qué idempotencia es necesaria
\end{itemize}

\section{Estructura sugerida}

\subsection*{Bloque 1: RPC (15 min)}

\begin{itemize}
\item presentar RPC
\item romper la analogía con función local
\item mostrar que una llamada remota no ocurre como un paso único
\end{itemize}

\subsection*{Bloque 1b: capas y no-inmediatez (15 min)}

\begin{itemize}
\item dibujar aplicación, cliente RPC, transporte, red, servidor y almacenamiento
\item preguntar qué observa cada capa ante pérdida o demora
\item construir una línea temporal donde el cliente abandona pero el servidor ejecuta
\end{itemize}

\subsection*{Bloque 2: timeout (15 min)}

\begin{itemize}
\item definir timeout
\item reforzar indistinguibilidad
\end{itemize}

\subsection*{Bloque 3: retry (20 min)}

\begin{itemize}
\item mostrar duplicación
\item construir timeline
\end{itemize}

\subsection*{Bloque 4: idempotencia (20 min)}

\begin{itemize}
\item definir
\item ejemplos
\item diseño
\end{itemize}

\subsection*{Bloque 5: integración (10 min)}

\begin{itemize}
\item aplicar modelo D
\end{itemize}

\section{Momento crítico}

Cuando el alumno entiende:

\begin{center}
\textbf{“retry es necesario, pero sin idempotencia rompe el sistema”}
\end{center}

\section{Errores frecuentes}

\begin{itemize}
\item creer que timeout = falla
\item pensar que la comunicación remota es instantánea
\item atribuir el problema a ``la red'' sin separar capas
\item creer que retry es seguro
\item ignorar duplicación
\end{itemize}

\section{Conexión con clase 3}

Esta clase lleva naturalmente a:

\begin{itemize}
\item consistencia
\item replicación
\item estado distribuido
\end{itemize}

\end{document}
